\section{What the method isolates}\label{sec:discussion}
Prime signs are independent, but their products satisfy exact identities
at arbitrarily large distances. The proof uses both facts. It counts the
homogeneous relation spaces created by multiplication, then uses the
prime environment either for an exact dependency graph or for an exact
one-word conditional coupling. In the latter route, potential dependence
is separated from the changes actually imposed by the forcing. The
three-term profile \eqref{eq:master-profile}, and its full-value extension
\eqref{eq:absolute-profile}, are the arithmetic outputs. The independent
Poisson field is obtained only after the separate marginal, local-overlap
and prime-support estimates have been verified.

\subsection{A class of observables, rather than one count}
The dictionary theorem separates local word combinatorics from nonlocal
arithmetic. Changing the prescribed signs does not change a valuation
matrix. Prescribing their absolute orientation removes at most two parity
conditions in a pair of windows; the host count controls the additional
relations. Consequently the same profile handles a growing dictionary and
retains the identity and position of every rare word. Explicit marker families make its overlap condition zero. For a typical
dictionary, one can do more: averaging before the worst-case estimate
reduces pair dependence to equality of the two complete words. The
finite collision bound removes the uniform cardinality restriction while
retaining the same labels and bounded total intensity. Short suffix--prefix matches can instead create
local clusters independently of any multiplicative relation; such
dictionaries require a separate analysis.

Long runs provide two complementary examples. Counting left-maximal starts
gives a Poisson field with geometric excesses and fair signs. Counting all
constant windows retains the clusters, and the field contracts to the
explicit compound-Poisson law of \cref{cor:window-clusters}. Both statements
come from the same comparison, not from two unrelated moment arguments.
The asymmetric almost-sure bounds in \cref{cor:as-longest} use its summable
quantitative errors; they do not require independence across scales.

The boundary is a different arithmetic regime. Its initial positive
cylinder imposes $\pi(L)$ independent prime constraints, whereas other
microscopic starts have an additional rank cost. The competition between
$2^{-\pi(L)}$ and $M2^{-L}$ therefore produces a change in the location
and sign of a rare long run. The two-clock theorem retains more than the
limiting source weights: it gives a prelimit lattice law, with prime-rank
overshoot at the boundary and integer overshoot in the bulk. The logistic
formula is a consequence of normalizing these two masses. The substantive
inputs are the microscopic rank surplus, the relative marked bulk kernel
and the negligible region between them.

\subsection{Arithmetic and probabilistic limits of the estimates}
Three terms reach $N^{5/3}$ at criticality in the present upper bound.
Improving only the terminal kernel energy would therefore not improve the
total exponent. The rational families in
\eqref{eq:rational-lower-dyadic}--\eqref{eq:rational-lower-macro} provide
lower exponents $4/3$ and $3/2$. Sharpening this gap is an arithmetic
question. It is distinct from improving the leading subexponential error,
which is governed by exceptional supports and shared prime coordinates.
For growing dictionaries the relevant quantity is not the full relation
sum, but its capped version. Mutual exclusivity at one site gives a
probability ceiling, and \cref{prop:capped-profile} preserves it through
the terminal count. This removes the terminal restriction on dictionary
size without improving the uncapped $5/3$ exponent. The remaining term
$NQ_B^{2/3}$ gives the sufficient exponent $1/2$ in
\eqref{eq:dictionary-critical}. Extending this uniform result requires control of the nonterminal
contribution or additional information about the dictionary. The typical
theorem provides the latter through its sampling law, not through a
smaller universal relation profile. A short affine description alone is
not a guarantee: the high-probability result concerns a uniformly sampled
full-rank matrix and an independent uniform right-hand side.

The Diophantine step uses only polynomial-height boxes. Localizing the
squarefree parts to the coefficient and root differences reduces the
count to two factors and a classical uniform Pell estimate; a height-free
bound on all integral points is not needed. The self-contained proof
fixes the required uniformity, rather than claiming priority for that
counting principle.

Different proof routes give sufficient budgets, not intrinsic prices
for retaining a label. In the hard corridor with $\Lambda\to\infty$,
\cref{thm:micro-run-tv,cor:dyadic-micro} retain all sites at the
one-factor budget, on a prefix and by restriction on a dyadic block.
Selected-prime reciprocals replace the all-large-prime footprint. The
tables in \cref{sec:conditional-laws} separate these available ranges
from the direct scalar and graph-based estimates. Local resolution still
requires an error negligible relative to the event being resolved.
At information-dependent cutoffs, \cref{prop:information-labelled}
balances deletion and graph costs; the event's measurability must be
retained along with its information cost.

\paragraph{What the Palm complement does and does not prove.}
\Cref{supp:sec:palm-complements} separates regular target configurations,
exact signed finite identities and consequences of the microscopic
comparison. Prescribed presence does not by itself control the absence
of further runs. The configuration-mass identity and
\cref{thm:micro-run-tv} imply the averaged one-sided Palm deficit
\cref{supp:palm:cor:palm-completion}; a stronger squared Walsh norm is
not asserted to vanish. Conversely,
\cref{supp:palm:thm:cumulant-obstruction} exhibits diverging all-order
absolute activity while microscopic distance tends to zero. Rare
configurations may dominate an exponential moment without preventing
ordinary total-variation approximation. The proof by selected pivots
never requires that activity to be small.

\subsection{Further questions}
Dictionaries with substantial short overlaps need a cluster analysis
compatible with the nonlocal arithmetic dependence. The compound-Poisson
word laws for Markov sources in
\cite{ReinertSchbath1998,RoquainSchbath2007} provide context, not an
automatic transfer. Our constant-window law is one model case, not a
classification of periodic or interacting patterns. Extending the
selected-coordinate coupling to nested word families, other prime
alphabets or larger conditioning fields would likewise require its own
exact forcing law, arithmetic footprint and tail estimates.

Records, censoring and occupation times require additional conventions
about the order and recognition of occurrences. Their process-level
analysis is outside the scope of this article. A common measurable
readout of the same-grid comparison does not by itself supply a new
Gaussian limit, a topology of convergence or asymptotic independence.
Similarly, the binary run-field theorem is not a comparison for every
prime-value distribution or for arbitrary persistent cluster targets.
These are boundaries of the results proved here, not assertions that
the corresponding questions are unresolved in all other models.
